% Preâmbulo LaTeX extra para o PDF do relatório (pandoc --include-in-header).
% Compilado com xelatex. Usa só pacotes presentes no TeX Live "basic".

% ---- Fontes ---------------------------------------------------------------
% Menlo tem glifos de box-drawing (│ ├ └ ─) da árvore de diretórios e cobre
% Greek/middot que aparecem dentro de blocos de código.
\usepackage{fontspec}
\setmonofont{Menlo}[Scale=0.82]

% ---- Imagens --------------------------------------------------------------
\usepackage{graphicx}
\setkeys{Gin}{width=0.85\linewidth,keepaspectratio}

% ---- Símbolos Unicode em texto corrido ------------------------------------
% O pacote newunicodechar não existe no BasicTeX; tornamos os caracteres ativos
% via kernel do XeTeX e os mapeamos para LaTeX robusto. Em blocos verbatim os
% catcodes são reiniciados, então dentro de código os caracteres caem na fonte
% mono (Menlo) — exatamente o desejado para a árvore de diretórios.
\catcode`→=\active \protected\def→{\ensuremath{\rightarrow}}
\catcode`↔=\active \protected\def↔{\ensuremath{\leftrightarrow}}
\catcode`α=\active \protected\def α{\ensuremath{\alpha}}
\catcode`β=\active \protected\def β{\ensuremath{\beta}}
\catcode`≈=\active \protected\def≈{\ensuremath{\approx}}
\catcode`≥=\active \protected\def≥{\ensuremath{\geq}}
\catcode`∈=\active \protected\def∈{\ensuremath{\in}}
\catcode`×=\active \protected\def×{\ensuremath{\times}}
\catcode`⁰=\active \protected\def⁰{\textsuperscript{0}}
\catcode`⁴=\active \protected\def⁴{\textsuperscript{4}}
\catcode`⁵=\active \protected\def⁵{\textsuperscript{5}}
\catcode`⁸=\active \protected\def⁸{\textsuperscript{8}}

% ---- Espaçamento ----------------------------------------------------------
\usepackage{parskip}
